\section{Benchmark construction}

\paragraph{Motivation from public chat logs.}
As a lightweight motivation check, we scan the first 20,000 streamed WildChat
conversations \citep{zhao2024wildchat} using predefined repair cues. Among
10,681 multi-turn conversations, 172 contain at least one follow-up cue related
to generic repair, preservation, output language, unwanted translation, script,
or register/locale. The scan records only aggregate counts and hashed metadata:
no raw user or assistant text is written. We use the scan to motivate the
failure taxonomy, not to estimate population prevalence.

\paragraph{Stress-pilot design.}
We created \textsc{RePromptTax-Stress-v0.2}, a 120-item synthetic benchmark
with a balanced $3\times4$ design: Spanish-English, Hindi-English, and
Arabic-English; and four task families with ten items per language-pair/family
cell. Table~\ref{tab:families} summarizes the families.

\begin{table}[t]
\centering
\small
\begin{tabular}{p{0.20\linewidth}p{0.28\linewidth}p{0.39\linewidth}}
\toprule
Family & User-side ambiguity & Expected contract \\
\midrule
Editing preservation & Instruction language differs from quoted content language. & Improve or rewrite the content while preserving its intended language; do not translate unless asked. \\
Output-language inference & Prompt mixes languages without an explicit ``answer in X'' instruction. & Infer the expected output language from the pragmatic context. \\
Quote preservation & The task concerns text around quoted headings or labels. & Preserve required quoted spans and literal data exactly. \\
Script/register/locale & The user specifies romanization, formal/informal register, or local literal formats. & Obey script, register, and locale constraints while preserving IDs, file names, times, and headings. \\
\bottomrule
\end{tabular}
\caption{The four task families in \textsc{RePromptTax-Stress-v0.2}. The pilot is designed to stress interaction contracts, not to estimate natural prevalence.}
\label{tab:families}
\end{table}

The benchmark is synthetic and privacy-safe. A release-hygiene audit confirms
120 unique prompts, zero normalized duplicate prompts, required scoring markers
for all rows, zero privacy-marker hits under the audit regexes, and 99 literal
preservation spans across the families that require them. The benchmark also
contains known-bad-output notes to clarify what each item is designed to catch.
These checks are not a substitute for native-speaker validation, but they make
the artifact easier to audit and reproduce.
